\documentclass[11pt]{article}
\usepackage[margin=1.2in]{geometry}
\usepackage[utf8]{inputenc}
\usepackage[T1]{fontenc}
\usepackage[english]{babel}
\usepackage{fourier}
\usepackage{amsthm}
\usepackage{amssymb}
\usepackage{amsmath}
\usepackage{amsfonts}
\usepackage{mathtools}
\usepackage{latexsym}
\usepackage{graphicx}
\usepackage{float}
\usepackage{etoolbox}
\usepackage{hyperref}
\usepackage{tikz}
\usepackage{booktabs}
\usepackage{pgfplots}
\pgfplotsset{compat=1.18} 
\usepackage{lipsum}
\usepackage{algorithm}
\usepackage{algpseudocode}
\usepackage{mathtools}
\usepackage{nccmath}
\usepackage[most]{tcolorbox}
\newtcolorbox[auto counter]{problem}[1][]{%
    enhanced,
    breakable,
    colback=white,
    colbacktitle=white,
    coltitle=black,
    fonttitle=\bfseries,
    boxrule=.6pt,
    titlerule=.2pt,
    toptitle=3pt,
    bottomtitle=3pt,
    title=GitHub repository of this project}

\usepackage{tensors}

\RequirePackage[activate={true,nocompatibility},final,tracking=true,kerning=true,spacing=true]{microtype}
\SetTracking{encoding={*}, shape=sc}{40} % Spacing for Small Caps


\newcommand{\K}{\mathbb{K}}
\newcommand{\R}{\mathbb{R}}
\newcommand{\N}{\mathbb{N}}
\newcommand{\Z}{\mathbb{Z}}
\newcommand{\Q}{\mathbb{Q}}
\newcommand{\C}{\mathbb{C}}
\newcommand{\V}{\mathcal{V}}
\newcommand{\E}{\mathcal{E}}
\newcommand{\G}{\mathcal{G}}
\renewcommand{\L}{\mathbb{L}}

\newcommand{\defeq}{\vcentcolon=}
\newcommand{\eqdef}{=\vcentcolon}
\newcommand{\norm}[1]{\left\lVert#1\right\rVert}
\renewcommand{\vec}[1]{\text{vec}#1}

\newtheorem{theorem}{Theorem}[section]
\newtheorem{lemma}[theorem]{Lemma}
\newtheorem{proposition}[theorem]{Proposition}
\newtheorem{corollary}[theorem]{Corollary}
\theoremstyle{definition}
\newtheorem{definition}[theorem]{Definition}
\newtheorem{example}[theorem]{Example}
\theoremstyle{remark}
\newtheorem{remark}[theorem]{Remark}

\title{%
    Accelerated Filtering on Graphs using Lanczos method
    \\ \large Scientific Computing project report}
\author{Alberto Defendi}

\date{}

\setlength{\parskip}{1em}
\setlength{\parindent}{0em}

\begin{document}

\maketitle


\begin{definition}
  Denotiamo l'insieme dei numeri naturali da $ 1 $ a $ n \in \N $ come $ [n] = \left\{ 1,\dots,n
  \right\} $, e il prodotto cartesiano tra $ [n] $ e $ [m] $ come $ [n] \otimes [m] =
  \left\{ (i,j) : i \in [n], j \in [m] \right\} $.
\end{definition}

\begin{definition}
  Un tensore $ \Tn{X}  $ di ordine $ d $ è un array multidimensionale con entrate definite sul campo $
  \K $, che denotiamo come $  \Tn{X} \in \Kmsiz{\K} $. Indicizziamo gli elementi di un tensore come
  $ \Tn{X}(\miwc) $, dove $ (\miwc) \in \mdom$.
\end{definition}

Analogamente alle matrici, i tensori vengono salvati in memoria come vettori, risulta quindi
naturale definire la seguente oprazione
\begin{definition}
  Dato un tensore $ \Tn{X}$, possiamo definire l'operazione che trasforma un tensore in vettore,
  tramite la mappa di vettorizzazione 
  $  \vec : \Rmsiz \to \R^{\prod_{j=1}^d n_j}$, dove $ x_l = \Tn{X}(\L^{-1}(l;\miwc)) $.
\end{definition}

\clearpage
\bibliographystyle{unsrt}
\bibliography{ref}
\nocite{*}



\end{document}
